\section{Las opciones en las que participas}
\begin{itemize}
  \item Artículo especializado para publicar en revista indizada
  \item Ensayo
  \item Memoria de experiencia laboral
  \item Reporte de aplicación de conocimientos
  \item Reporte de autoempleo profesional
  \item Reporte de residencia de investigación
  \item Tesina
  \item Tesis
\end{itemize}
\section{Lo que es igual en todas las opciones}
Las pantallas de esta sección se tomaron en la opción \emph{Tesis}; en las demás opciones son iguales.
\subsection{Firma del voto aprobatorio}
Abre el enlace, consulta el documento u oficio y firma (\boton{Firmar} o \boton{Firmar de recibido}).
\captura{tesis_sinodal_01_firma_del_voto_aprobatorio}{Enlace personal del correo «Firma del voto aprobatorio»: lees el documento y pulsas «Firmar».}
\subsection{Oficio de citación}
Abre el enlace, consulta el documento u oficio y firma (\boton{Firmar} o \boton{Firmar de recibido}).
\captura{tesis_sinodal_02_oficio_de_citacion}{Enlace personal «Oficio de citación»: consultas tu oficio y firmas de recibido.}
\section{Lo propio de cada opción}
\subsection{Artículo especializado para publicar en revista indizada}
Trabajo escrito de esta opción: Artículo original e inédito, enviado a una revista inscrita en índices validados. Como revisor firmas el voto aprobatorio; como sinodal, el acuse de tu oficio para el examen previo y para el definitivo.
\subsection{Ensayo}
Trabajo escrito de esta opción: Exposición crítica y propositiva de un tema, con postura intelectual propia. Como revisor firmas el voto aprobatorio; como sinodal, el acuse de tu oficio para el examen previo y para el definitivo.
\subsection{Memoria de experiencia laboral}
Trabajo escrito de esta opción: Recopilación y resultados de una experiencia profesional en un ámbito reconocido. Como revisor firmas el voto aprobatorio; como sinodal, el acuse de tu oficio para el examen previo y para el definitivo.
\subsection{Reporte de aplicación de conocimientos}
Trabajo escrito de esta opción: Propuesta de solución a un problema determinado, derivada de lo aprendido en la carrera. Como revisor firmas el voto aprobatorio; como sinodal, el acuse de tu oficio para el examen previo y para el definitivo.
\subsection{Reporte de autoempleo profesional}
Trabajo escrito de esta opción: Informe sobre la constitución y desarrollo de una unidad económica formal. Como revisor firmas el voto aprobatorio; como sinodal, el acuse de tu oficio para el examen previo y para el definitivo.
\subsection{Reporte de residencia de investigación}
Trabajo escrito de esta opción: Informe de actividades de apoyo a la investigación bajo la dirección de un investigador activo en el SNI. Como revisor firmas el voto aprobatorio; como sinodal, el acuse de tu oficio para el examen previo y para el definitivo.
\subsection{Tesina}
Trabajo escrito de esta opción: Investigación documental de carácter monográfico, con aportación al estado del arte. Como revisor firmas el voto aprobatorio; como sinodal, el acuse de tu oficio para el examen previo y para el definitivo.
\subsection{Tesis}
Trabajo escrito de esta opción: Investigación documental, experimental, empírica o teórica; individual o de hasta tres pasantes. Como revisor firmas el voto aprobatorio; como sinodal, el acuse de tu oficio para el examen previo y para el definitivo.
